\scanpage{83}

% L1
espaces forment un drapeau.

% L2–6
e) Pour que $G$ opérant sur $C$ soit irréductible, il f et s. que $\rho$ et les $\beta_i$ ($0\leq i\leq n-1$) soient $\ne0$, i.e. que \deleted{\unsure{$f$}} la forme \deleted{quadratique can. $f_V$} \inserted{bil. can. $\unsure{\varphi_V|\psi_V|f_V}$} sur $V$ soit \deleted{non dég.} non dégénérée (i.e. que $f_V$ soit lisse, et que \deleted{\unsure{pour}} \inserted{si} $n$ impair, \deleted{$2$ soit inv. en $S$} \inserted{$k$ soit de car. $\ne2$})

% L7
\deleted{\unsure{f)|c)} Pour que}

% L8–11
\uline{Dém} Dire que $Z$ est inv. par $\unsure{\tau_i|\mathrm T_i}$ signifie que l'on a $Z\ni h_i$ ou $Z\subset H_i$. \deleted{Il est} \inserted{Supposons le} invariant par tous les $\tau_i$, \deleted{ou bien $Z\subset H_i$ pour tt $0\leq i\leq n$, donc $Z\subset\bigcap H_i=\{c\}$, donc $Z=\varnothing$ ou $Z=\{c\}$, ou bien $\exists\,i$}

% L11–14
Notons que si $0<i\leq\unsure{n|n+1}$ $Z\subset H_i\Longrightarrow Z\subset H_{i-1}$, car sinon on aurait $h_{i-1}\in\unsure{Z|H_i}$ donc $h_{i-1}\notin H_i$, absurde. \unsure{De même|D'ailleurs}, si $h_j\in Z$, \deleted{\unsure{$\forall j$|$h_j$}} $0\leq j\unsure{<n|\leq n}$, alors $h_{j+1}\in Z$, car sinon on aurait $Z\subset H_{j+1}$ donc $h_j\notin H_{j+1}$, absurde. \deleted{\unsure{D}}Il en résulte \deleted{\unsure{qu'il}} \inserted{qu'on} a deux possibilités

% L17–18
a) $Z\subset H_i\qquad \forall\ 0\leq i\leq n$, donc $Z\subset\bigcap_{0\leq i\leq n}H_i=\{c\}$, donc $Z=\varnothing$, ou $Z=\{c\}$ -- qui sont bien invariants

% L19–21
a$'$) $Z\not\subset H_i\qquad \forall\ 0\leq i\leq n$ donc $Z\ni h_i\quad 0\leq\forall i\leq n$ donc $Z\supset\unsure{V|\vee}(h_0\,\ldots\,h_n)=C$, donc $Z=C$ ou \inserted{$Z=$}$X$ qui sont bien invariants

% L22–23
b) On n'est pas en cas a) ni b), donc $\exists\ i$, $0\leq i<n$, tel que $Z\subset H_j$ \inserted{si $0\leq j\leq i$} et $Z\not\subset H_j$ si $n\geq j\geq i+1$

% archive sheet number: 82
